\chapter{Persistencia: redb e índices secundarios}
\label{cap:persistencia}

El cap.~\ref{cap:implementacion} mostró cómo el TypeGraph vive en
heap. Este capítulo cubre cómo se vuelca a disco y se recupera, y por
qué el truco crítico de la persistencia no es la serialización sino la
\emph{rehidratación de los índices secundarios}.

\section{El backend redb}

\citaCrate{sotfs-storage} es un wrapper minimalista (~85 LOC) sobre
\redb{}, una base de datos KV embebida con transacciones ACID. Hay una
sola tabla:

\begin{lstlisting}[language=rust]
const METADATA_TABLE: TableDefinition<&str, &[u8]>
    = TableDefinition::new("metadata");

pub struct RedbBackend { db: Database, }
\end{lstlisting}

La tabla guarda un único par: \texttt{("graph", <bytes>)}. Los bytes
son el \texttt{TypeGraph} entero serializado con \texttt{serde\_json}.
La elección es deliberadamente simple: la atomicidad la da la
transacción de \redb, no un protocolo casero.

\section{Save y load}

\paragraph{Save.} (\citaLine{sotfs-storage/src/backend.rs}{29})

\begin{algorithm}[H]
\SetAlgoLined
\KwIn{Backend $B$, grafo $g$}
\caption{\textsc{Backend.Save}}\label{alg:save}

$\textit{bytes} \gets \texttt{serde\_json::to\_vec}(g)$\;
$\textit{txn} \gets B.\textrm{db}.\texttt{begin\_write}()$\;
$\textit{table} \gets \textit{txn}.\texttt{open\_table}(\texttt{METADATA\_TABLE})$\;
$\textit{table}.\texttt{insert}(\texttt{"graph"}, \textit{bytes})$\;
$\textit{txn}.\texttt{commit}()$
\end{algorithm}

\paragraph{Load.} (\citaLine{sotfs-storage/src/backend.rs}{41})

\begin{algorithm}[H]
\SetAlgoLined
\KwIn{Backend $B$}
\KwOut{Grafo $g$ rehidratado, o \texttt{None}}
\caption{\textsc{Backend.Load}}\label{alg:load}

$\textit{txn} \gets B.\textrm{db}.\texttt{begin\_read}()$\;
$\textit{table} \gets \textit{txn}.\texttt{open\_table}(\texttt{METADATA\_TABLE})$\;
\If{$\textit{table}.\texttt{get}(\texttt{"graph"}) = \texttt{None}$}{
  \Return{\texttt{None}}\;
}
$g \gets \texttt{serde\_json::from\_slice}(\textit{bytes})$\;
$g.\textsc{RebuildDirNameIdx}()$
\Return{\texttt{Some}(g)}\;
\end{algorithm}

\section{El truco: rehidratación de índices secundarios}

El \texttt{TypeGraph} contiene varios índices que no se serializan:

\begin{itemize}
  \item \texttt{dir\_name\_idx: BTreeMap<(DirId, String), EdgeId>} —
    Marcado \texttt{\#[serde(skip)]}.
  \item \texttt{dir\_for\_inode\_idx: BTreeMap<InodeId, DirId>} —
    \texttt{\#[serde(default)]}, vacío tras deserialización.
\end{itemize}

\paragraph{¿Por qué saltarlos?} Las claves tupla \texttt{(DirId,
String)} no caben naturalmente en JSON. Más importante: estos índices
son \emph{datos derivados}, no estado canónico. Pueden reconstruirse a
partir de \texttt{dir\_contains} (que sí se serializa). Persistirlos
duplicaría datos y agregaría riesgo de inconsistencia.

\paragraph{La función crítica.} \texttt{rebuild\_dir\_name\_idx}
(\citaLine{sotfs-graph/src/graph.rs}{673}) recorre cada arista
\econtains{} viva y la inserta en el \texttt{BTreeMap}:

\begin{algorithm}[H]
\SetAlgoLined
\KwIn{Grafo $g$ recién deserializado}
\caption{\textsc{Graph.RebuildDirNameIdx}}\label{alg:rebuild-idx}

$g.\textrm{dir\_name\_idx} \gets \emptyset$\;
\ForEach{$(d, \textit{eids}) \in g.\textrm{dir\_contains}$}{
  \ForEach{$\textit{eid} \in \textit{eids}$}{
    \If{$g.\textsc{GetEdge}(\textit{eid}) = \texttt{Contains}\{name, \dots\}$}{
      $g.\textrm{dir\_name\_idx}.\textsc{Insert}((d, name), \textit{eid})$\;
    }
  }
}
\end{algorithm}

\paragraph{Costo.} $O(M \log M)$ donde $M$ es el número de aristas
\econtains. Para un FS de \(10^5\) entradas, sub-segundo.

\paragraph{¿Por qué crítico?} Sin la rehidratación, \texttt{lookup\_name}
caería al fallback lineal $O(N)$
(cap.~\ref{cap:typegraph} \trampref{tr:on-cube}); el filesystem
funcionaría pero degradaría a 92\% de CPU bajo workload realista.

\section{El comando \texttt{sotfsctl}}

\citaCrate{sotfs-cli} expone:

\begin{description}
  \item[\texttt{sotfsctl mkfs <db.redb>}] crea un FS vacío con root inode
    + root dir.
  \item[\texttt{sotfsctl check <db.redb>}] carga el FS, ejecuta
    \texttt{check\_invariants()} y \texttt{check\_dir\_name\_idx\_consistency()}.
    Versión \emph{offline} del fsck.
  \item[\texttt{sotfsctl dump <db.redb>}] vuelca a stdout el grafo en
    DOT o JSON.
\end{description}

\section{Crash safety: lo que da redb}

\begin{teorema}[Garantías de redb]
\label{teo:redb-garantias}
La transacción de \redb{} provee atomicidad, durabilidad e isolation
serializable. Concretamente: si \texttt{txn.commit()} retorna
\texttt{Ok}, los datos están en NAND tras el power-cut; si retorna
\texttt{Err} o crashea durante el commit, el estado pre-transacción
sobrevive intacto.
\end{teorema}

La consecuencia para sotFS: el modelo formal del cap.~\ref{cap:transacciones}
con su \wal{} explícito no se implementa en el runtime. La capa de
durabilidad la provee \redb. Si \redb{} cumple sus contratos, sotFS
hereda crash safety.

\begin{trampa}[fsync sin barrier en SSDs consumer-grade]
\label{tr:fsync-ssd}
\redb{} llama a \texttt{fsync(2)} antes de retornar de \texttt{commit}.
Pero un SSD consumer-grade que ignora el comando \texttt{FUA} puede
ackear el \texttt{fsync} antes de flushear su caché DRAM volátil. Bajo
power-cut, los bytes ``persistidos'' se pierden.

\textbf{Mitigación.} El operador del FS debe verificar el SSD bajo
power-cut físico (típicamente con un test de \emph{torn write}). Modelo
asume \fsync{} honesto. Discusión completa en cap.~\ref{cap:transacciones}.
\end{trampa}

\section*{Resumen del capítulo}

\begin{itemize}
  \item Backend \redb: una sola tabla, un solo blob JSON, atomicidad
    via \texttt{txn.commit}.
  \item Save serializa el TypeGraph entero; load deserializa más
    \emph{rehidrata} los índices secundarios via
    \texttt{rebuild\_dir\_name\_idx}.
  \item La rehidratación es \(O(M \log M)\) y crítica: sin ella,
    \texttt{lookup\_name} cae a $O(N)$ y degrada todo el FS.
  \item \texttt{sotfsctl mkfs/check/dump} expone las primitivas para
    inicializar, verificar y volcar el FS offline.
  \item Crash safety se hereda de las garantías de \redb; el modelo TLA
    con \wal{} explícito vive un nivel de abstracción más arriba.
\end{itemize}

\section*{Ejercicios}

\begin{ejercicio}[\dificultad{$\star$}]
\label{ej:pers-1}
¿Qué tabla(s) usa \citaCrate{sotfs-storage}? ¿Qué claves y valores?
\end{ejercicio}

\begin{ejercicio}[\dificultad{$\star$}]
\label{ej:pers-2}
¿Por qué \texttt{dir\_name\_idx} no se serializa con
\texttt{serde\_json}?
\end{ejercicio}

\begin{ejercicio}[\dificultad{$\star\star$}]
\label{ej:pers-3}
Estime el costo de save / load para un FS de \(10^5\) inodes.
Considerá serialización JSON de un \texttt{TypeGraph} de ese tamaño.
\end{ejercicio}

\begin{ejercicio}[\dificultad{$\star\star$}]
\label{ej:pers-4}
Diseñe un test que verifique la consistencia del índice
\texttt{dir\_name\_idx} tras un round-trip save/load + 1000 ops más.
\end{ejercicio}

\begin{ejercicio}[\dificultad{$\star\star\star$}]
\label{ej:pers-5}
Reemplace el blob JSON por un schema de varias tablas redb: una para
inodes, una para edges, etc. Discuta tradeoffs en velocidad de
save/load, fragmentación on-disk, y atomicidad cross-table.
\end{ejercicio}
